\documentclass[11pt]{article}

\usepackage[margin=1in]{geometry}
\usepackage[utf8]{inputenc}
\usepackage[T1]{fontenc}
\usepackage{amsmath,amssymb,amsthm}
\usepackage{graphicx}
\usepackage{booktabs}
\usepackage{array}
\usepackage{longtable}
\usepackage{xcolor}
\usepackage{hyperref}
\usepackage{enumitem}
\usepackage{caption}
\usepackage{float}
\usepackage{parskip}

\hypersetup{
  colorlinks=true,
  linkcolor=blue!50!black,
  citecolor=blue!50!black,
  urlcolor=blue!50!black,
  pdftitle={The BAE Bertini Project},
  pdfauthor={Zuxian He}
}

\newtheorem{definition}{Definition}[section]
\newtheorem{proposition}{Proposition}[section]
\newtheorem{corollary}{Corollary}[section]
\newtheorem{remark}{Remark}[section]
\newtheorem{conjecture}{Empirical Symmetry}[section]
\newtheorem{example}{Example}[section]

\newcommand{\file}[1]{\texttt{#1}}
\newcommand{\code}[1]{\texttt{#1}}
\newcommand{\Wron}{\operatorname{Wr}}
\DeclareMathOperator{\Res}{Res}

\title{\textbf{Wronskian Solutions of (Supersymmetric) Bethe Ansatz Equations:\\
Formalism and a Bertini-Based Numerical Continuation Pipeline}}
\author{Zuxian He \\ \small Department of Physics, Uppsala University}
\date{September 2026}

\begin{document}

\maketitle

\begin{abstract}
We give a formal account of the mathematical construction and numerical
realization implemented in the \file{BAE\_Bertini} codebase, which computes
Bethe roots of a nested integrable spin chain for eigenstates labeled by
standard Young tableaux (SYT). The construction proceeds through a
\emph{Q-system}: a lattice of auxiliary polynomials $Q_{a,s}(u)$ obeying a
discrete bilinear (Hirota/QQ) functional relation, solved not directly but
by representing the $Q_{a,s}$ as Wronskian (Casoratian) determinants of a
common set of elementary building blocks --- a discrete analogue of the
classical fact that Wronskians of solutions of a shared linear equation
obey simple conservation laws. We derive the degree counting of
the Q-functions from the geometry of the Young diagram, present the
supersymmetric (bosonic/fermionic) generalization of the Wronskian
solution obtained by choosing a \emph{marked point} (equivalently a
Kac--Dynkin grading) on the diagram, and derive the vacancy-number
combinatorics of the associated rigged-configuration / string-hypothesis
description. We then formalize the numerical strategy: a single chain
inhomogeneity is treated as a deformation parameter $\Lambda$, and the
Q-system is solved by a \emph{parameter homotopy continuation} from an
asymptotically free ($\Lambda\to\infty$) point, where the solution is
combinatorially exact, down to the physical, homogeneous point
$\Lambda=0$. We recall the general theory of homotopy continuation (the
Davidenko differential equation and its predictor-corrector discretization)
and describe how the project delegates this step to the numerical
algebraic geometry package \textbf{Bertini}, through an adaptive
multiple-precision path tracker, invoked via a parameter-homotopy driver
written in Python. We close with the combinatorial symmetry used to
validate numerical solutions (a rigging-complementation involution,
realized combinatorially as Sch\"utzenberger evacuation, under which Bethe
roots are conjectured to negate), quantitative validation and performance
results, and a discussion of the construction's scope and limitations.
\end{abstract}

\tableofcontents
\newpage

%=============================================================================
\section{Introduction}
%=============================================================================

\subsection{Physical setting}

Integrable spin chains of rank higher than one admit a \emph{nested}
Bethe ansatz: magnon excitations are organized into several coupled
families of rapidities, one family per node of an auxiliary root system,
and the full spectrum is obtained by solving a coupled system of Bethe
ansatz equations (BAE) across all families simultaneously. For a chain of
length $L$, the spectrum organizes into multiplets labeled by partitions
$\lambda$ of $L$ (Young diagrams); within a multiplet of shape $\lambda$,
individual eigenstates are labeled by the standard Young tableaux (SYT)
of that shape --- fillings of the diagram's $L$ boxes with the integers
$1,\dots,L$ that increase along every row and down every column. This
report is a self-contained account of a computational pipeline that,
given a single SYT, numerically constructs the corresponding Bethe roots.

Rather than attacking the Bethe equations directly in root variables, the
pipeline uses the modern \emph{Q-system} reformulation of the nested Bethe
ansatz, in which the fundamental objects are polynomials $Q_{a,s}(u)$ of
the spectral parameter $u$, attached to nodes $(a,s)$ of a lattice built
from the Young diagram, and required to satisfy a discrete bilinear
functional relation (Section~\ref{sec:qsystem}). The distinguishing
feature of this project is that the Q-system is not solved by treating
the coefficients of every $Q_{a,s}$ as independent unknowns; instead, a
closed-form solution is constructed directly as a \textbf{Wronskian
(Casoratian) determinant} of a small number of elementary building
blocks, split into a \emph{bosonic} and a \emph{fermionic} family
according to a chosen \emph{marked point} on the diagram
(Section~\ref{sec:susy}). This is the origin of the name
\emph{(SUSY) Wronskian Bethe Equations} carried by the project.

\subsection{Numerical strategy}

Even after the Wronskian ansatz reduces the problem to finitely many
polynomial equations in finitely many coefficients, those equations must
still be solved numerically for each SYT. The pipeline does this by
introducing a single auxiliary parameter $\Lambda$ --- realized physically
as one chain site's inhomogeneity --- and constructing, at
$\Lambda\to\infty$, an exact combinatorial solution associated with the
tableau's shape (Section~\ref{sec:asymptotics}). The physical, homogeneous
point $\Lambda=0$ is then reached by \emph{numerically continuing} this
solution as $\Lambda$ is driven from its large starting value to zero ---
a textbook parameter homotopy continuation problem, which we formalize in
Section~\ref{sec:homotopy}. The tableau itself is built up one box at a
time, each new box's site briefly appearing at large $\Lambda$ and then
being continued into the homogeneous chain, so that a single SYT
computation is in fact a sequence of $|\lambda|$ homotopy continuations,
each seeded by the previous one's converged solution
(Section~\ref{sec:sequential}).

The actual path tracking is delegated to \textbf{Bertini}, a numerical
algebraic geometry package implementing adaptive multiple-precision
predictor-corrector continuation, through its Python interface. The
symbolic construction of the Wronskian relations, and the orchestration of
many such continuations across all SYT of a Young diagram, is carried out
in the Wolfram Language; the two layers communicate through an explicit
polynomial system exported to a plain-text file. We describe this
numerical layer formally in Section~\ref{sec:homotopy} and summarize its
software realization in Section~\ref{sec:implementation}.

\subsection{Scope and evidentiary status of this report}

This report is based on a direct reading of the project's source files. A
number of internal comments cite specific equation numbers of an external
reference (``Eq.~2.3'', ``Eq.~2.33'', ``Eq.~5.7'', ``Eq.~5.8''), indicating
that the construction of Section~\ref{sec:susy} is a direct numerical
implementation of a specific paper; that paper is not included in the
repository. Where this report states a physical interpretation that is not
explicit in an in-code comment, this is flagged as an inference from the
code's structure rather than a fact quoted from a reference. The
mathematical structure itself --- the Q-system, the Wronskian solution,
the marked-point/grading freedom, the vacancy-number combinatorics, and
the numerical continuation scheme --- is, by contrast, directly and
unambiguously present in the source and is derived formally below.

\subsection{Outline}

Section~\ref{sec:qsystem} introduces the Q-system and explains the
Wronskian mechanism behind the Hirota/QQ relation. Section~\ref{sec:susy}
constructs the supersymmetric Wronskian solution and its marked-point
freedom. Section~\ref{sec:asymptotics} introduces the deformation
parameter $\Lambda$, its asymptotic solution, and the sequential
construction. Section~\ref{sec:rc} develops the rigged-configuration /
string-hypothesis combinatorics and the flip symmetry used for
validation. Section~\ref{sec:homotopy} formalizes homotopy continuation
and Bertini's realization of it. Section~\ref{sec:worked-example} works
every preceding definition out by hand on the smallest nontrivial shape.
Section~\ref{sec:implementation} summarizes the software pipeline.
Section~\ref{sec:results} reports validation and performance results.
Section~\ref{sec:discussion} discusses scope, limitations, and outlook.
Table~\ref{tab:notation} below collects the recurring notation for
reference.

\begin{table}[H]
\centering
\begin{tabular}{@{}ll@{}}
\toprule
\textbf{Symbol} & \textbf{Meaning} \\
\midrule
$\lambda\vdash L$, $\lambda'$ & Young diagram ($L$ boxes) labeling a multiplet, and its conjugate \\
$\mathrm{SYT}$ & standard Young tableau: a filling of $\lambda$ labeling one eigenstate \\
$u$ & spectral parameter (rapidity) \\
$Q_{a,s}(u)$ & Q-function at lattice node $(a,s)$; $M(a,s)$ its polynomial degree \\
$\theta_1,\dots,\theta_L$ & chain-site inhomogeneities; all zero for a homogeneous chain \\
$\Wron[f,g]$ & discrete Wronskian (Casoratian) of $f,g$, Eq.~\eqref{eq:wron2} \\
$\mathrm{mP}=(\lambda_\ast,\nu_\ast)$ & marked point: bosonic/fermionic split of elementary Q-functions \\
$Q_{A|J}$, $\gamma_{A,J,k}$ & Q-function (resp.\ its coefficients) for bosonic set $A$, fermionic set $J$ \\
$c_{a,s,k}$ & generic-basis Taylor coefficients of $Q_{a,s}$ (Section~\ref{sec:bases}) \\
$\Lambda$, $\Lambda_0$ & continuation (deformation) parameter and its large starting value \\
$n^{(a)}_s$, $J$, $P^{(a)}_s$ & string multiplicity, rigging, and vacancy number (Section~\ref{sec:rc}) \\
$t$, $H(x,t)$ & homotopy parameter and homotopy function (Section~\ref{sec:homotopy}) \\
\bottomrule
\end{tabular}
\caption{Recurring notation.}
\label{tab:notation}
\end{table}


%=============================================================================
\section{The Q-System Formalism}
\label{sec:qsystem}
%=============================================================================

Written directly in terms of Bethe roots, the equations of a nested Bethe
ansatz are a system of coupled transcendental equations in which every
root of every family scatters off every other root of every family: even
before any numerics, simply \emph{writing down} the equations for a
chain with several nesting levels and more than a handful of roots per
level is already unwieldy, and finding a good starting guess for a
numerical solver is far from obvious. The Q-system reformulation
sidesteps both difficulties at once. Rather than working with the roots
themselves, one works with the (much better-behaved) generating
polynomials that have those roots as their zeros; the coupling between
different families of roots is then encoded not in a system of rational
equations but in a single, uniform bilinear \emph{functional} relation
between neighboring polynomials, Eq.~\eqref{eq:qq} below, which is the
same at every node of the lattice regardless of how many roots are
involved. This uniformity is exactly what makes the construction
amenable to the closed-form (Wronskian) solution method of the rest of
this section, and to the systematic numerical continuation of
Section~\ref{sec:homotopy}.

\subsection{Q-functions on the Young diagram}

Let $\lambda=(\lambda_1\ge\lambda_2\ge\cdots)$ be a partition of $L$, drawn
as a Young diagram, and let $\lambda'$ denote its conjugate partition. To
each lattice point $(a,s)$ with $0\le a$, $0\le s$ we associate a monic
polynomial $Q_{a,s}(u)$ of the spectral parameter $u$, of degree
\begin{equation}
  \label{eq:degree}
  M(a,s) \;=\; L + a\,s \;-\; \sum_{i=1}^{a}\lambda_i \;-\; \sum_{j=1}^{s}\lambda'_j ,
\end{equation}
so that $M(a,s)=0$ once $(a,s)$ leaves the diagram's ``co-hook'' (the
region where the formula would otherwise become negative); such nodes
carry the trivial function $Q_{a,s}\equiv1$. The distinguished node
$(0,0)$ has $M(0,0)=L$ and its polynomial,
\begin{equation}
  \label{eq:vacuum}
  Q_{0,0}(u) \;=\; \prod_{n=1}^{L}\bigl(u-\theta_n\bigr),
\end{equation}
is fixed entirely by the chain's $L$ inhomogeneities $\theta_1,\dots,\theta_L$
(one per site); for a homogeneous chain, all $\theta_n=0$ except possibly a
single deformed site, as in Section~\ref{sec:asymptotics}.

\begin{example}[A running example: $\lambda=(2,1)$]
\label{ex:shape21}
Throughout this report we will repeatedly return to the smallest Young
diagram with more than one row, $\lambda=(2,1)$: two rows, of lengths $2$
and $1$, three boxes in total ($L=3$). Its conjugate is $\lambda'=(2,1)$
as well (the diagram is self-conjugate: two boxes stacked in the first
column, one box in the second). This shape has exactly two standard
fillings,
\[
  T_A=\begin{smallmatrix}1&2\\3&\end{smallmatrix}, \qquad
  T_B=\begin{smallmatrix}1&3\\2&\end{smallmatrix},
\]
matching the hook-length count $3!/(3\cdot1\cdot1)=2$; in the language of
the symmetric group $S_3$ acting on three objects, $\lambda=(2,1)$ is the
two-dimensional ``standard'' irreducible representation, and $T_A,T_B$
label a basis of it. Evaluating the degree formula \eqref{eq:degree} at
every node with $a+s\le 2$ gives
\[
  M(0,0)=3, \qquad M(1,0)=M(0,1)=3-2=1, \qquad M(1,1)=M(2,0)=M(0,2)=0,
\]
so only three nodes carry a nontrivial polynomial: the length-3 vacuum
$Q_{0,0}$, and two length-1 (linear) polynomials $Q_{1,0}$, $Q_{0,1}$ ---
each with a single root. Since $\lambda$ has two rows, there is exactly
one nesting level (Eq.~\eqref{eq:bethe-roots} below), so this shape's
Bethe ansatz has, per eigenstate, a single Bethe root: the zero of
$Q_{1,0}$.
\end{example}

\subsection{The Hirota (QQ) equation}
\label{sec:hirota}

Write $f^{[\pm k]}(u):=f(u\pm\tfrac{ik}{2})$ for a shift of the spectral
parameter by $k$ half-units along the imaginary axis. The Q-functions of
the nested Bethe ansatz are required to satisfy, at every node $(a,s)$ for
which all four neighbors are defined,
\begin{equation}
  \label{eq:qq}
  Q_{a,s}^{[+1]}\,Q_{a+1,s+1}^{[-1]} \;-\; Q_{a,s}^{[-1]}\,Q_{a+1,s+1}^{[+1]}
  \;=\; i\bigl(M(a,s)-M(a+1,s+1)\bigr)\, Q_{a+1,s}(u)\, Q_{a,s+1}(u).
\end{equation}
This is the (additive-normalization) \emph{Hirota bilinear equation}, also
called the \emph{QQ-relation}, of the discrete integrable hierarchy
underlying the nested Bethe ansatz. It is, on the face of it, a
transcendental functional equation; the project's numerical strategy
depends on converting it into finitely many \emph{polynomial} equations in
finitely many unknowns, which is accomplished in two complementary ways
described below.

\subsubsection*{Polynomial form}

If every $Q_{a,s}$ appearing in \eqref{eq:qq} is taken to be a polynomial
of the degree prescribed by \eqref{eq:degree}, with undetermined Taylor
coefficients $c_{a,s,k}$,
\begin{equation}
  Q_{a,s}(u) \;=\; u^{M(a,s)} + \sum_{k=1}^{M(a,s)} c_{a,s,k}\, u^{M(a,s)-k},
\end{equation}
then the left- and right-hand sides of \eqref{eq:qq} are themselves
polynomials in $u$ (the half-integer shifts only ever produce integer
powers of $i$ once the two half-shifted products are subtracted), and
matching coefficients of each power of $u$ turns \eqref{eq:qq} into a
finite list of \emph{polynomial equations} in the $c_{a,s,k}$. This
``generic'' route is available in the codebase and defines the ambient
polynomial system that any solution --- however constructed --- must
satisfy.

\subsubsection*{Wronskian (Casoratian) solution}

The project's principal method does not solve the polynomial system of
the previous paragraph directly. Instead it \emph{generates} a family of
functions satisfying \eqref{eq:qq} identically, by representing each
$Q_{a,s}$ as a determinant of shifted copies of a common set of
``elementary'' functions --- a discrete analogue of the Wronskian
determinant of solutions of a linear ODE, sometimes called a
\emph{Casoratian}. Concretely, define the discrete Wronskian of two
functions $f,g$ by
\begin{equation}
  \label{eq:wron2}
  \Wron[f,g](u) \;:=\; f^{[+1]}(u)\,g^{[-1]}(u) \;-\; f^{[-1]}(u)\,g^{[+1]}(u).
\end{equation}

Why should a determinant of this kind have any chance of satisfying a
bilinear relation like \eqref{eq:qq}? The intuition is a discrete analogue
of a classical fact about ordinary Wronskians: if two functions $f,g$ both
solve a common linear second-order difference (``Baxter-type'') equation
with shared potential $V(u)$,
\begin{equation}
  \label{eq:baxter-toy}
  X(u+i) + X(u-i) \;=\; V(u)\,X(u), \qquad X\in\{f,g\},
\end{equation}
then their discrete Wronskian is a \emph{conserved quantity} of the shift:
\begin{align}
  \label{eq:wronskian-conservation}
  \Wron[f,g]^{[+1]}(u) \;-\; \Wron[f,g]^{[-1]}(u)
  &\;=\; g(u)\bigl[f(u{+}i){+}f(u{-}i)\bigr] \;-\; f(u)\bigl[g(u{+}i){+}g(u{-}i)\bigr] \notag\\
  &\;=\; V(u)\bigl[g(u)f(u)-f(u)g(u)\bigr] \;=\;0,
\end{align}
by direct substitution of \eqref{eq:baxter-toy} for both $f$ and $g$ ---
exactly the discrete counterpart of the elementary fact that the
Wronskian of two solutions of a linear ODE with no first-derivative term
is constant. This is the generic phenomenon exploited throughout the
Q-system literature: elementary Q-functions are required to be solutions
of a shared linear difference equation fixed by the vacuum data
\eqref{eq:vacuum}, and Casoratians of such solutions automatically satisfy
conservation-law-type bilinear identities under the shift $u\to u\pm i$.

\begin{remark}[What is, and is not, claimed here]
Equation~\eqref{eq:wronskian-conservation} is only the simplest instance
of this mechanism (a vanishing right-hand side, for two solutions of the
\emph{same}, ungraded linear equation). The actual QQ-relation
\eqref{eq:qq} that the project's Q-system must satisfy has a
\emph{nonzero} right-hand side $i(M_{a,s}-M_{a+1,s+1})\,Q_{a+1,s}Q_{a,s+1}$,
and its bosonic and fermionic elementary building blocks
(Section~\ref{sec:susy}) are solutions of two \emph{different, graded}
linear problems, not one shared ungraded equation. Establishing that the
specific super-Wronskian ansatz used here satisfies every instance of
\eqref{eq:qq} on the full lattice --- not merely the toy identity
\eqref{eq:wronskian-conservation} --- is a genuine consistency theorem
about that graded construction; it belongs to the reference work the
implementation follows (cited there by equation number, as noted in the
Introduction) and is not re-derived from first principles in this report.
What \emph{is} directly confirmed in the source code, and is recorded
below exactly as implemented, is the recursive definition by which every
higher Q-function is \emph{generated} from elementary ones once that
consistency is assumed.
\end{remark}

The project's implementation makes this recursion explicit and constructive:
writing $Q_{A|J}$ for the Q-function attached to a bosonic index set $A$
and a fermionic index set $J$ (Section~\ref{sec:susy}), one has the
reduction
\begin{align}
  \label{eq:reduction-bos}
  Q_{A_0\cup\{a,b\}\,|\,J} &\;=\; \frac{\Wron\bigl[\,Q_{A_0\cup\{a\}|J},\,Q_{A_0\cup\{b\}|J}\,\bigr]}{Q_{A_0|J}}, \\
  \label{eq:reduction-ferm}
  Q_{A\,|\,J_0\cup\{i,j\}} &\;=\; \frac{\Wron\bigl[\,Q_{A|J_0\cup\{i\}},\,Q_{A|J_0\cup\{j\}}\,\bigr]}{Q_{A|J_0}},
\end{align}
i.e.\ a Q-function with one more bosonic (resp.\ fermionic) index than
$Q_{A_0|J}$ is the ratio of the Wronskian of its two ``parents'' to
$Q_{A_0|J}$ itself --- the discrete analogue of a classical Plücker
relation among Wronskians of nested subsets of solutions. Every
multi-index Q-function of the system is generated from single-index
($|A|=|J|=1$, or one of $A,J$ empty) elementary Q-functions by iterating
\eqref{eq:reduction-bos}--\eqref{eq:reduction-ferm}.


%=============================================================================
\section{The Supersymmetric Wronskian and the Marked-Point Freedom}
\label{sec:susy}
%=============================================================================

\subsection{Marked points and gradings}

The recursion \eqref{eq:reduction-bos}--\eqref{eq:reduction-ferm} treats
the elementary building blocks as falling into two families, indexed by
disjoint label sets: a \emph{bosonic} family of size $\lambda_\ast$ and a
\emph{fermionic} family of size $\nu_\ast$, with $\lambda_\ast+\nu_\ast$
equal to the number of rows of the Young diagram participating in the
construction. The project calls the choice
$\mathrm{mP}=(\lambda_\ast,\nu_\ast)$ a \textbf{marked point}: it fixes how
many of the diagram's rows are assigned to the bosonic family versus the
fermionic one. In the representation theory of a superalgebra
$\mathfrak{gl}(\lambda_\ast|\nu_\ast)$, this is precisely the datum of a
choice of \emph{simple-root system} (a \textbf{Kac--Dynkin diagram}): unlike
an ordinary Lie algebra, a Lie superalgebra admits several inequivalent
but representation-theoretically equivalent Dynkin diagrams, related by
``fermionic duality'' moves, and different choices lead to Wronskian
solutions of different total polynomial degree. The implementation makes
this freedom explicit by enumerating admissible marked points/paths across
the diagram and selecting one that minimizes the total number of Bethe
roots to be computed --- a computational-efficiency criterion, not a
physical one, applied on top of the mathematical freedom just described.

\begin{remark}
The repository does not state explicitly whether the physical spin chain
being solved is itself supersymmetric, or whether the bosonic/fermionic
split is used purely as a computational device that produces an efficient
closed-form Wronskian solution for an ordinary (non-super) nested Bethe
ansatz problem. Both readings are consistent with the mathematics above;
this report presents the construction in a manner agnostic to that
distinction and defers to the external source cited internally by
equation number for the precise physical statement.
\end{remark}

\begin{example}[Marked points for $\lambda=(2,1)$]
\label{ex:markedpoints21}
Continuing Example~\ref{ex:shape21}: the diagram has two rows, so a
marked point is a splitting $2=\lambda_\ast+\nu_\ast$, and there are three
choices, $(\lambda_\ast,\nu_\ast)\in\{(2,0),(1,1),(0,2)\}$. The choice
$(2,0)$ assigns \emph{both} rows to the bosonic family and none to the
fermionic one: the super-Cauchy determinant \eqref{eq:susywronskian}
below degenerates to an ordinary $2\times2$ Wronskian determinant of two
purely bosonic towers, with no fermionic block at all --- the ``purely
bosonic'' grading. The choice $(1,1)$ instead assigns one row to each
family, mixing a single bosonic and a single fermionic elementary
Q-function; the choice $(0,2)$ assigns both rows to the fermionic family.
All three choices describe the \emph{same} vacuum polynomial $Q_{0,0}$, by
construction, once each is separately fit to the target inhomogeneities
\eqref{eq:wronskian-relations} --- they differ only in how many elementary
coefficients must be solved for along the way, which is exactly the
efficiency trade-off the marked-point search of
Section~\ref{sec:implementation} is built to exploit.
\end{example}

\subsection{The super-Wronskian determinant}

Fix a marked point $\mathrm{mP}=(\lambda_\ast,\nu_\ast)$. Denote the
$\lambda_\ast$ bosonic elementary Q-functions by $Q_{\{p\}|\emptyset}$,
$p=1,\dots,\lambda_\ast$, and the $\nu_\ast$ fermionic ones by
$Q_{\emptyset|\{q\}}$, $q=1,\dots,\nu_\ast$. Form the two rectangular
blocks of shifted elementary functions
\begin{align}
  B_{p,j}(u) &= Q_{\{p\}|\emptyset}\Bigl(u+\tfrac{i}{2}(\lambda_\ast-\nu_\ast-2j+1)\Bigr),
  && 1\le p\le\lambda_\ast,\ \ 1\le j\le \lambda_\ast-\nu_\ast, \\
  F_{p,q}(u) &= Q_{\{p\}|\{q\}}\Bigl(u+\tfrac{i}{2}(\lambda_\ast-\nu_\ast)\Bigr),
  && 1\le p\le\lambda_\ast,\ \ 1\le q\le \nu_\ast,
\end{align}
where the mixed bosonic--fermionic building block $Q_{\{p\}|\{q\}}$ is
itself fixed, for each single pair $(p,q)$, by demanding that it solve the
elementary instance of \eqref{eq:qq} coupling the purely bosonic and
purely fermionic elementary Q-functions,
\begin{equation}
  Q_{\{p\}|\{q\}}^{[+1]} - Q_{\{p\}|\{q\}}^{[-1]} \;=\; Q_{\{p\}|\emptyset}\; Q_{\emptyset|\{q\}},
\end{equation}
solved order by order in $u$ for the coefficients of $Q_{\{p\}|\{q\}}$
(this is a linear recursion in the Taylor coefficients, since the
right-hand side is already known once the bosonic and fermionic
elementary blocks are fixed). The vacuum Q-function is then the single
$\lambda_\ast\times\lambda_\ast$ super-Cauchy/Wronskian determinant
\begin{equation}
  \label{eq:susywronskian}
  Q_{0,0}(u) \;=\; (-1)^{\nu_\ast(\lambda_\ast-\nu_\ast)}\,
  \det\bigl[\, F \mid B \,\bigr](u),
\end{equation}
i.e.\ the determinant of the $\lambda_\ast\times\lambda_\ast$ matrix whose
first $\nu_\ast$ columns are the fermionic block $F$ and whose remaining
$\lambda_\ast-\nu_\ast$ columns are the bosonic block $B$. Built up by
iterating the Wronskian-quotient recursion
\eqref{eq:reduction-bos}--\eqref{eq:reduction-ferm} from the elementary
bosonic and fermionic blocks, \eqref{eq:susywronskian} automatically
satisfies the QQ-relation \eqref{eq:qq} against its neighbors for
\emph{any} choice of the elementary block coefficients (Remark following
Eq.~\eqref{eq:wronskian-conservation}); the only remaining task is to fix
those coefficients so that \eqref{eq:susywronskian} reproduces the
physical vacuum polynomial \eqref{eq:vacuum}, i.e.\ so that
\begin{equation}
  \label{eq:wronskian-relations}
  \operatorname{monic}\bigl[\,Q_{0,0}^{\text{Wronskian}}(u)\,\bigr] \;=\; \prod_{n=1}^{L}(u-\theta_n),
\end{equation}
which, after clearing the overall normalization and matching coefficients
of each power of $u$, is a finite polynomial system in the elementary
blocks' coefficients --- the \textbf{Wronskian relations} that are
ultimately handed to the numerical continuation of
Section~\ref{sec:homotopy}.

\subsection{Two coefficient bases}
\label{sec:bases}

Because the physical content of the Q-system is carried by the elementary
bosonic/fermionic building blocks, while the generic polynomial system of
Section~\ref{sec:hirota} is expressed in the coefficients $c_{a,s,k}$ of
\emph{every} node, the implementation maintains an explicit two-way
translation between:
\begin{itemize}
  \item the \textbf{generic basis}, $\{c_{a,s,k}\}$, one set of Taylor
  coefficients per Young-diagram node;
  \item the \textbf{Wronskian basis}, $\{\gamma_{A,J,k}\}$, the Taylor
  coefficients of the elementary bosonic/fermionic blocks introduced
  above.
\end{itemize}
The translation is carried out by a \emph{triangular elimination}: at each
stage, one selects a coefficient-matching equation from
\eqref{eq:wronskian-relations} (or its analogue for other nodes) that
contains exactly one still-undetermined variable while every other symbol
appearing in it is already known, and solves for that variable; the
process is iterated until either every unknown has been eliminated or no
such equation remains (in which case the system is under-determined in
that basis, a known limitation for certain small diagrams). This is the
formal content of what the implementation calls, informally, the
``strict triangular sort''; only the Wronskian-basis coefficients
$\gamma_{A,J,k}$ are exported as unknowns to the numerical continuation
step, and the generic basis is recovered afterward, purely as a
convenience for extracting Bethe roots (Section~\ref{sec:root-extraction}
below).


%=============================================================================
\section{The Deformation Parameter and the Sequential Construction}
\label{sec:asymptotics}
%=============================================================================

\subsection{The chain inhomogeneity as a homotopy parameter}
\label{sec:lambda-def}

Consider building the tableau's chain one site at a time. After the
$(n{-}1)$-th site has already been solved (all its Wronskian-basis
coefficients fixed at the physical, homogeneous point), the $n$-th site is
introduced by setting its inhomogeneity to a formal parameter $\Lambda$
while every earlier site retains its already-converged, homogeneous value
of $0$:
\begin{equation}
  \label{eq:vacuum-lambda}
  Q_{0,0}^{(n)}(u;\Lambda) \;=\; u^{\,n-1}\,(u-\Lambda).
\end{equation}
The physical (homogeneous) chain of length $n$ is recovered at
$\Lambda=0$; the value $\Lambda\to\infty$ instead decouples the new site
entirely, in the following precise sense.

\subsection{Asymptotic solution at large $\Lambda$}

Every root of every Q-function in the Wronskian construction can be
tracked, order by order in $1/\Lambda$, by attaching to \emph{each root} a
formal weight built from the tableau's combinatorics: writing $\theta_m$
for a single formal deformation parameter attached to tableau entry
$m\in\{1,\dots,L\}$, one sets
\begin{equation}
  \label{eq:theta-rep}
  q_{a,s,k} \;=\; \Bigl(\prod \text{hook-content ratios of the sub-tableau at }(a,s,k)\Bigr)\,\theta_{m(a,s,k)},
\end{equation}
a combinatorial (rational-function-of-diagram-data) weight multiplying the
single formal parameter of whichever tableau entry that root is
associated with. Specializing $\theta_m=\Lambda$ for the newest entry
($m=n$) and $\theta_m=0$ for every earlier entry, and substituting into
the Wronskian relations \eqref{eq:wronskian-relations} together with
\eqref{eq:vacuum-lambda}, yields a system that is exactly solvable ---
typically linear or of very low degree --- because the bulk of the
tableau's roots are pinned to $\theta_m=0$ and only the newest entry's
roots carry any $\Lambda$-dependence at all. This closed-form (or
near-closed-form) solution, valid asymptotically as $\Lambda\to\infty$,
supplies the starting point for the numerical continuation: no root-finding
is required to seed it, only the combinatorial evaluation of
\eqref{eq:theta-rep} on the given tableau.

\subsection{Sequential construction}
\label{sec:sequential}

The full computation for a target SYT of shape $\lambda\vdash L$
proceeds inductively on the chain of sub-tableaux obtained by restricting
the target tableau to entries $\{1\},\{1,2\},\dots,\{1,\dots,L\}$ (the
tableau's own growth sequence, a saturated chain
$\emptyset=\lambda^{(0)}\subset\lambda^{(1)}\subset\cdots\subset\lambda^{(L)}=\lambda$
in Young's lattice):
\begin{enumerate}
  \item at $n=1$, the one-site chain is solved by the exact asymptotic
  construction of the previous subsection, continued (trivially) to
  $\Lambda=0$;
  \item at each subsequent $n$, the $(n{-}1)$-site converged solution is
  held fixed, the new site's inhomogeneity is set to a symbolic $\Lambda$
  as in \eqref{eq:vacuum-lambda}, the exact large-$\Lambda$ solution
  \eqref{eq:theta-rep} is evaluated as a numerical starting point at a
  large finite value of $\Lambda$, and the Wronskian relations are then
  numerically continued in $\Lambda$ from that large value down to
  $\Lambda=0$ (Section~\ref{sec:homotopy}), yielding the converged
  $n$-site solution.
\end{enumerate}
Different standard Young tableaux of the same shape correspond to
different orderings of this box-by-box growth, hence to different
continuation paths; they are expected, on general representation-theoretic
grounds, to terminate on the distinct members of the degenerate
eigenspace associated with that shape. The total computation for one SYT
is therefore not a single homotopy continuation but a sequence of $L-1$
of them (the first site requiring none), each a full instance of the
formalism developed in Section~\ref{sec:homotopy}, seeded by the previous
one's output.

\begin{example}[Growth sequences for $T_A$ and $T_B$]
Continuing Example~\ref{ex:shape21}, the two tableaux of shape $(2,1)$
trace out \emph{different} intermediate shapes on their way to
$\lambda=(2,1)$, even though they end at the same shape:
\[
  T_A=\begin{smallmatrix}1&2\\3&\end{smallmatrix}:
  \quad \emptyset \to (1) \to (2) \to (2,1),
  \qquad\qquad
  T_B=\begin{smallmatrix}1&3\\2&\end{smallmatrix}:
  \quad \emptyset \to (1) \to (1,1) \to (2,1).
\]
Both start identically (a single site, trivially continued) and end
identically (three sites, shape $(2,1)$), but their \emph{second} step
continues from a one-site chain to a genuinely different two-site
sub-problem: $T_A$'s second site is added to the first row (an
intermediate shape $(2)$, with no second row yet, hence a trivial,
un-nested Q-system at that step), while $T_B$'s second site opens the
second row (an intermediate shape $(1,1)$, already carrying one nesting
level). This concretely illustrates why the two tableaux, despite sharing
a final shape and hence a final set of Wronskian relations
\eqref{eq:wronskian-relations}, are continued along different paths and
are expected to converge on the two distinct eigenstates that shape
$(2,1)$'s degenerate multiplet contains.
\end{example}

\subsection{Extraction of Bethe roots}
\label{sec:root-extraction}

Once the coefficients of every node have been recovered in the generic
basis $\{c_{a,s,k}\}$ (Section~\ref{sec:bases}) at the physical point
$\Lambda=0$ for the full tableau, the Bethe roots themselves are the
zeros of the \emph{first-column} Q-functions,
\begin{equation}
  \label{eq:bethe-roots}
  Q_{w,0}(u) \;=\; 0, \qquad w=1,\dots,(\text{number of rows of }\lambda)-1,
\end{equation}
one independent root set per nesting level $w$, consistent with a nested
Bethe ansatz of rank equal to the number of rows of the diagram.


%=============================================================================
\section{Rigged Configurations, the String Hypothesis, and a Flip Symmetry}
\label{sec:rc}
%=============================================================================

\subsection{Strings and vacancy numbers}

Independently of the Wronskian construction above, the classical
\emph{string hypothesis} conjectures that solutions of the (ordinary,
non-super) nested Bethe equations organize into \textbf{strings}: for each
node $a$ of the auxiliary root system, a collection of complex roots
sharing a common real center $v$ and equally spaced along the imaginary
direction,
\begin{equation}
  \label{eq:string}
  u_k \;=\; v - \frac{i}{2}(s-1-2k), \qquad k=0,1,\dots,s-1,
\end{equation}
for a string of length $s$. A configuration of strings is encoded
combinatorially by a \textbf{rigged configuration}: for each node $a$ and
each string length $s$, a nonnegative integer multiplicity
$n^{(a)}_s$ together with, for every one of those $n^{(a)}_s$ strings, an
integer \textbf{rigging} $J\in\{0,1,\dots,P^{(a)}_s\}$ bounded by the
node's \textbf{vacancy number}
\begin{equation}
  \label{eq:vacancy}
  P^{(a)}_s \;=\; \delta_{a,1}\,L \;-\; \sum_{b} A_{ab}\sum_{s'} \min(s,s')\, n^{(b)}_{s'},
\end{equation}
where $A_{ab}$ is the Cartan matrix of the auxiliary (type-$A_{n-1}$, in
the cases implemented here) root system. Equation~\eqref{eq:vacancy} is
the standard counting-function/vacancy formula of the string hypothesis:
$\delta_{a,1}L$ is the leading contribution of the physical chain to the
first node, and the Cartan-matrix contraction against the ``string-string
scattering kernel'' $\min(s,s')$ subtracts, for every other string present
at neighboring nodes, the number of momentum slots it excludes at node
$a$. A rigging $J$ for a given string is then, in the usual
interpretation, a choice of one of the $P^{(a)}_s+1$ admissible mode
numbers left available after this exclusion --- exactly analogous to
filling single-particle levels subject to a generalized Pauli principle.
The number of admissible riggings for a fixed string-length content is
therefore
\begin{equation}
  \label{eq:kostka-count}
  \prod_{a,s} \binom{P^{(a)}_s + n^{(a)}_s}{n^{(a)}_s},
\end{equation}
a product of binomial coefficients that (summed appropriately over all
admissible string contents of a given node) reproduces the Kostka-number
counting of standard Young tableaux --- the combinatorial statement of
completeness of the Bethe ansatz.

\begin{example}[Vacancy numbers for $\lambda=(2,1)$]
\label{ex:vacancy21}
The shape $(2,1)$ of Example~\ref{ex:shape21} has a single nesting level
(Eq.~\eqref{eq:bethe-roots}), so its auxiliary root system has a single
node $a=1$, with Cartan entry $A_{11}=2$. Consider the string content with
one string of length $s=1$ and nothing else, $n^{(1)}_1=1$: by
\eqref{eq:vacancy},
\[
  P^{(1)}_1 \;=\; \delta_{1,1}\cdot 3 \;-\; A_{11}\cdot\min(1,1)\cdot n^{(1)}_1
  \;=\; 3 - 2\cdot1\cdot1 \;=\; 1,
\]
so the single string's rigging can be $J=0$ or $J=1$: exactly
$\binom{P+n}{n}=\binom{1+1}{1}=2$ admissible riggings, matching the two
standard Young tableaux $T_A,T_B$ of that shape one for one. This is the
smallest nontrivial instance of the general fact
\eqref{eq:kostka-count}: the rigged-configuration count of admissible
riggings for a fixed string content reproduces the number of standard
Young tableaux of the corresponding shape.
\end{example}

\subsection{The Kerov--Kirillov--Reshetikhin bijection}

The \textbf{Kerov--Kirillov--Reshetikhin (KKR) bijection} is a classical,
constructive correspondence between rigged configurations and standard
Young tableaux (equivalently, crystal paths in a tensor product of
single-box representations). The project implements both directions of
this bijection explicitly, following the classical algorithm: recursively
remove (or add) the box associated with the current largest (or next)
tableau entry, choosing at each step the unique \emph{singular} string ---
one whose rigging equals its vacancy number exactly --- as the string a
box is added to or removed from, and updating every vacancy number via
\eqref{eq:vacancy} after each such move. The construction converts a
rigged configuration into explicit \textbf{mode numbers} for each string
(via a further shift of the rigging by the vacancy number and the string
index) which, through the string ansatz \eqref{eq:string}, give an
explicit approximate (string-hypothesis) prediction for the Bethe roots
associated with a tableau, independent of the exact Wronskian
construction of Sections~\ref{sec:qsystem}--\ref{sec:susy}.

\subsection{A flip involution and a conjectured root symmetry}
\label{sec:flip}

Complementing every string's rigging against its own vacancy number,
\begin{equation}
  \label{eq:flip}
  J \;\longmapsto\; P^{(a)}_s - J,
\end{equation}
applied simultaneously to every string of a rigged configuration, defines
an involution on rigged configurations (and, via the KKR bijection, an
involution on standard Young tableaux of a fixed shape), which we denote
$\mathrm{SYT}\mapsto\mathrm{Flip}(\mathrm{SYT})$.

\begin{conjecture}[Root negation under the flip]
\label{conj:flip}
For a standard Young tableau $\mathrm{SYT}$ with Bethe roots
$\{u_1,\dots,u_r\}$ (at any fixed nesting level), the flipped tableau's
Bethe roots satisfy
\begin{equation}
  \{u_1,\dots,u_r\} \;=\; \{-u_1,\dots,-u_r\}
  \quad\text{as computed for }\mathrm{Flip}(\mathrm{SYT}).
\end{equation}
\end{conjecture}

This statement is not derived in the codebase from first principles; it
is used purely as an \emph{empirical, numerically checked} consistency
condition (Section~\ref{sec:validation-flip}). Combinatorially, the flip
\eqref{eq:flip} coincides, on standard Young tableaux, with
\textbf{Sch\"utzenberger evacuation} --- the classical involution obtained
by repeatedly removing the box containing the tableau's smallest entry via
jeu-de-taquin slides and relabeling --- a fact from crystal-base theory
(evacuation realizes the crystal-graph analogue of complex conjugation of
the underlying representation). The project exploits this coincidence
purely as a fast, dependency-free way to compute the flipped tableau
without constructing its rigged configuration explicitly, and treats
Conjecture~\ref{conj:flip} as an independent numerical cross-check on the
Wronskian-continuation output of Sections~\ref{sec:qsystem}--\ref{sec:asymptotics}
rather than as a derivation of it.


%=============================================================================
\section{Numerical Continuation: Formalism and Realization}
\label{sec:homotopy}
%=============================================================================

\subsection{A concrete warm-up example}
\label{sec:warmup}

Before setting up the general formalism, it is worth fixing ideas with the
simplest possible instance, reproduced directly from the project's own
introductory material. Take the target system
\begin{equation}
  \label{eq:toy-target}
  F(x,y) \;=\; \begin{pmatrix} x^2+y^2-1 \\ x+y \end{pmatrix} \;=\; 0,
\end{equation}
i.e.\ the intersection of the unit circle with the line $y=-x$, with
exact solutions $(x,y)=(\pm\tfrac{1}{\sqrt2},\mp\tfrac{1}{\sqrt2})$. This
system has no free parameter to continue in, so one instead introduces an
artificial \emph{start system} $G(x,y)=0$ of the same degrees ($2$ and
$1$, respectively) with known solutions --- a \textbf{total-degree start
system} --- and a homotopy $H=(1-t)F+tG$ exactly as in
\eqref{eq:homotopy-general} below. Bézout's theorem bounds the number of
solution paths that must be tracked by the product of the two degrees,
$2\times1=2$; correspondingly there are two start points to track, one
per path, from $t=1$ down to $t=0$. Tracking both paths with an adaptive
multiple-precision corrector (as described below) recovers the two exact
intersection points to the tracker's tolerance, with real parts
$\approx\pm0.70710678\ldots=\pm1/\sqrt2$ --- this is exactly what one
finds running Bertini on \eqref{eq:toy-target}, and it is the same
mechanism, applied to a genuine polynomial system with no natural
deformation parameter of its own, that Section~\ref{sec:parameter-homotopy}
specializes to the project's actual problem, where instead of an
artificial start system one already has an \emph{exact} solution at
$\Lambda=\Lambda_0$ to continue from.

\subsection{Homotopy continuation and the Davidenko equation}

Let $F:\mathbb{C}^N\to\mathbb{C}^N$ be the target polynomial system (here,
the Wronskian relations \eqref{eq:wronskian-relations} at $\Lambda=0$,
in the coefficients $\gamma$ of Section~\ref{sec:bases}), and let $G$ be a
\emph{start system} with known solution(s), continuously connected to $F$
by a homotopy
\begin{equation}
  \label{eq:homotopy-general}
  H(x,t) \;=\; (1-t)\,F(x) + t\,G(x), \qquad t\in[0,1],
\end{equation}
so that $H(x,1)=G(x)$ and $H(x,0)=F(x)$. Assuming the solution path
$x(t)$ defined implicitly by $H(x(t),t)=0$ is smooth, differentiating with
respect to $t$ gives the \textbf{Davidenko differential equation},
\begin{equation}
  \label{eq:davidenko}
  J_x H\bigl(x(t),t\bigr)\, \dot x(t) \;+\; \partial_t H\bigl(x(t),t\bigr) \;=\; 0
  \quad\Longrightarrow\quad
  \dot x(t) \;=\; -\bigl[J_x H(x(t),t)\bigr]^{-1}\,\partial_t H\bigl(x(t),t\bigr),
\end{equation}
an ordinary differential equation for the path $x(t)$, provided the
Jacobian $J_x H$ remains nonsingular along the path. Numerically tracking
the path from $t=1$ to $t=0$ then reduces to (i) a \textbf{predictor} step
that numerically integrates \eqref{eq:davidenko} by one of the standard
explicit ODE schemes, and (ii) a \textbf{corrector} step that removes the
resulting drift off the true solution path by Newton iteration,
\begin{equation}
  \label{eq:newton}
  x_{k+1} \;=\; x_k \;-\; \bigl[J_x H(x_k,t)\bigr]^{-1} H(x_k,t),
\end{equation}
applied at fixed $t$ until convergence to a specified tolerance. This
predictor--corrector scheme, together with \emph{adaptive} control of both
the step size in $t$ and the working numerical precision (raised whenever
the Jacobian in \eqref{eq:newton} becomes ill-conditioned, as happens near
a path singularity), is exactly what Bertini's \textbf{adaptive
multiple-precision tracker} (``AMPTracker'') implements; the several
named predictors available (constant, Euler, Heun, and a family of
embedded Runge--Kutta schemes of increasing order) are simply different
choices of explicit ODE integrator for \eqref{eq:davidenko}, trading
predictor accuracy (and hence achievable step size) against per-step
cost.

\subsection{Specialization: a linear parameter homotopy in $\Lambda$}
\label{sec:parameter-homotopy}

The construction of Section~\ref{sec:asymptotics} is a particularly
favorable instance of \eqref{eq:homotopy-general}: rather than requiring
an artificial start system, the Wronskian relations
\eqref{eq:wronskian-relations}, left symbolic in $\Lambda$, already
constitute a one-parameter polynomial family $F_\Lambda(\gamma)=0$ whose
member at a large value $\Lambda=\Lambda_0$ has the exact combinatorial
solution of Section~\ref{sec:asymptotics}, and whose member at $\Lambda=0$
is the physical system. Introducing the affine reparametrization
\begin{equation}
  \label{eq:param-path}
  \Lambda(t) \;=\; (1-t)\cdot 0 + t\cdot\Lambda_0 \;=\; t\,\Lambda_0, \qquad t\in[0,1],
\end{equation}
gives $H(\gamma,t):=F_{\Lambda(t)}(\gamma)$, with $H(\gamma,1)=F_{\Lambda_0}(\gamma)$
solved exactly and $H(\gamma,0)=F_0(\gamma)$ the target; exactly one path
is tracked (there is no Bézout-number multiplicity of start points to
enumerate, since the start point is already known exactly), from $t=1$ to
$t=0$. This is the parameter homotopy actually solved by the numerical
back end: the polynomial system exported at each step of the sequential
construction (Section~\ref{sec:sequential}) is $F_\Lambda(\gamma)$ with
$\Lambda$ left as the single homotopy variable and \eqref{eq:param-path}
substituted in symbolically before the system is handed to the tracker,
together with the exact starting value $\gamma(\Lambda_0)$ from
Section~\ref{sec:asymptotics} as the $t=1$ start point.

\subsection{Numerical safeguards}

Two safeguards are imposed around the tracked endpoint before it is
accepted as a physical solution. First, since the physical Bethe roots
recovered downstream (Eq.~\eqref{eq:bethe-roots}) are real linear
combinations of the tracked coefficients for the tableaux considered here,
each returned coefficient $\gamma=\gamma_{\mathrm{re}}+i\gamma_{\mathrm{im}}$
is required to satisfy
\begin{equation}
  \label{eq:real-check}
  |\gamma_{\mathrm{im}}| \;\le\; \varepsilon_{\mathrm{track}}\cdot\max(1,|\gamma_{\mathrm{re}}|),
\end{equation}
for the tracker's own tolerance $\varepsilon_{\mathrm{track}}$, with any
violation (or any non-finite value) aborting the run rather than
silently truncating an unexpectedly complex endpoint. Second, every
numerical result is published to disk only by an atomic replace of a
completed temporary file, so that a reader can never observe a partially
written solution, and a failed or timed-out continuation simply leaves no
new output rather than a corrupted one.

\subsection{A path-tracking failure mode: branch jumps}
\label{sec:branch-jump-formal}

The predictor--corrector scheme of \eqref{eq:davidenko}--\eqref{eq:newton}
is only guaranteed to follow the correct solution branch if the step size
in $t$ is small enough, relative to the local curvature of the path and
the separation of nearby solution branches, for the corrector to converge
back to the \emph{same} branch the predictor was following. When this
fails --- typically because the path passes close to another solution
branch or to a singularity of the discriminant locus --- the tracker can
converge to a numerically valid but physically incorrect solution: a
\emph{branch jump}. Empirically, for a pair of tableaux related by the
flip involution of Section~\ref{sec:flip} (whose roots are conjectured to
be exact negatives of one another, Conjecture~\ref{conj:flip}), starting
the continuation from a substantially larger $\Lambda_0$ and tracking at
higher precision and tighter tolerance was found to eliminate an observed
branch-jump discrepancy entirely (Section~\ref{sec:validation-flip}); both
adjustments reduce, respectively, the total path length that must be
tracked and the size of the correction the Newton step must make at each
point along it, and so plausibly reduce the risk of such a jump, though
neither is a formal guarantee against it in general.

\begin{remark}[Why extra precision helps]
It is worth making the mechanism a little more concrete. Two solution
branches that pass close to one another correspond to the Jacobian
$J_xH(x(t),t)$ in \eqref{eq:newton} becoming nearly singular somewhere
along the path (formally, the path approaches the \emph{discriminant
locus} where two roots of the system coincide). Near such a point, the
Newton correction $[J_xH]^{-1}H$ is the product of a huge, ill-conditioned
inverse and a small residual, so the correction's accuracy is extremely
sensitive to round-off in $H$ and $J_xH$ themselves: at fixed (say,
machine) precision, the correction step can easily be thrown to the
\emph{wrong} one of the two nearby branches. Raising the working
precision shrinks the round-off in $H$ and $J_xH$ directly, which is
exactly what lets the corrector resolve which of the two nearby branches
it started closest to --- this is the sense in which \emph{adaptive}
precision (raised locally, only where the path actually passes near such
a point, rather than uniformly everywhere) is doing real, not merely
cosmetic, numerical work.
\end{remark}


%=============================================================================
\section{A Fully Worked Example: the Shape \texorpdfstring{$(2,1)$}{(2,1)}}
\label{sec:worked-example}
%=============================================================================

It is worth pausing to see every piece introduced so far --- the degree
count, the number of nesting levels, the flip involution, and the
expected root negation --- meet in a single, completely explicit,
hand-checkable calculation, before turning to the software that automates
it for shapes too large to solve by hand.

For the purely bosonic marked point $(\lambda_\ast,\nu_\ast)=(2,0)$ of
Example~\ref{ex:markedpoints21}, the shape $\lambda=(2,1)$ reduces to the ordinary (non-super)
nested Bethe ansatz for a homogeneous chain of $L=3$ sites, with one
nesting level and, per Example~\ref{ex:shape21}, a single Bethe root per
eigenstate. This is precisely the textbook \emph{one-magnon sector} of the
periodic Heisenberg (XXX) spin chain: a state with a single magnon at
rapidity $u$ is an eigenstate of the transfer matrix provided the
periodicity condition
\begin{equation}
  \label{eq:one-magnon}
  \left(\frac{u+\tfrac{i}{2}}{u-\tfrac{i}{2}}\right)^{\!L} \;=\; 1
\end{equation}
holds --- the $M=1$ case of the Bethe equations, in which the usual
scattering factors between magnons are simply absent. At $L=3$,
clearing denominators in \eqref{eq:one-magnon} gives
$(u+\tfrac{i}{2})^3=(u-\tfrac{i}{2})^3$, and since
$a^3-b^3=(a-b)(a^2+ab+b^2)$ with $a-b=i\ne0$, this reduces to the single
quadratic
\begin{equation}
  \label{eq:one-magnon-quadratic}
  \Bigl(u+\tfrac{i}{2}\Bigr)^{\!2} + \Bigl(u+\tfrac{i}{2}\Bigr)\Bigl(u-\tfrac{i}{2}\Bigr) + \Bigl(u-\tfrac{i}{2}\Bigr)^{\!2}
  \;=\; 3u^2 - \tfrac14 \;=\; 0,
  \qquad\Longrightarrow\qquad
  u \;=\; \pm\frac{1}{2\sqrt3}.
\end{equation}
There are exactly two roots, matching the two standard Young tableaux
$T_A,T_B$ of Example~\ref{ex:shape21} one for one: each tableau's
degree-one Q-function $Q_{1,0}(u)=u\mp\tfrac{1}{2\sqrt3}$ contributes one
of the two roots of \eqref{eq:one-magnon-quadratic}, consistent with the
general multiplicity count of Section~\ref{sec:rc}
(two admissible riggings, Example~\ref{ex:vacancy21}, for two SYT).

A direct jeu-de-taquin computation (the evacuation algorithm of
Section~\ref{sec:flip}) confirms that $T_A$ and $T_B$ are each other's
flip, $\operatorname{Flip}(T_A)=T_B$: removing the box containing entry
$1$ from $T_A=\{\{1,2\},\{3\}\}$, sliding the resulting hole down and
relabeling, reproduces $T_B=\{\{1,3\},\{2\}\}$ exactly. Conjecture~\ref{conj:flip}
therefore predicts $\operatorname{Roots}(T_A)=-\operatorname{Roots}(T_B)$
--- and indeed the two roots found in
\eqref{eq:one-magnon-quadratic}, $+\tfrac{1}{2\sqrt3}$ and
$-\tfrac{1}{2\sqrt3}$, are exact negatives of one another. This is the
smallest possible instance in which every structural claim of this
report --- the degree/multiplicity count, the two-tableaux labeling of a
degenerate multiplet, and the flip-negation symmetry --- can be checked
by hand, with no numerical continuation required at all (the one-magnon
equation \eqref{eq:one-magnon} is exactly solvable). For any shape with
more than one magnon per level, no closed form is available in general,
and this is exactly where the asymptotic construction and homotopy
continuation of Sections~\ref{sec:asymptotics}--\ref{sec:homotopy} take
over.


%=============================================================================
\section{Software Realization}
\label{sec:implementation}
%=============================================================================

This section summarizes, at the level of architecture rather than
implementation detail, how the formalism of Sections~\ref{sec:qsystem}--\ref{sec:homotopy}
is realized as a working pipeline; it is included for completeness and to
ground the validation results of Section~\ref{sec:results}, not as a
substitute for the mathematical description above.

\textbf{Symbolic layer (Wolfram Language).} Given a Young diagram, a
symbolic layer constructs the Q-system lattice, selects a marked point,
builds the super-Wronskian expression \eqref{eq:susywronskian}, and
derives the Wronskian relations \eqref{eq:wronskian-relations} together
with the two-basis translation of Section~\ref{sec:bases}. A second,
independent module implements the asymptotic construction of
Section~\ref{sec:asymptotics} and drives the sequential, box-by-box
computation of Section~\ref{sec:sequential} for a given standard Young
tableau, delegating each individual continuation step to a pluggable
numerical back end.

\textbf{Numerical back end (Bertini, via Python).} At each continuation
step, the current Wronskian relations (with $\Lambda$ left symbolic per
\eqref{eq:param-path}) and a high-precision starting point are exported to
a plain-text table (one polynomial equation per row, each row carrying
the variable name, its starting value, and its defining expression). A
Python driver parses this table into a genuine polynomial system, forms
the parameter path of \eqref{eq:param-path}, and tracks a single path with
Bertini's adaptive multiple-precision tracker (Section~\ref{sec:homotopy}),
exposing the tracker's tolerance, precision, step-size, Newton-iteration,
and predictor-scheme parameters as configurable settings. The converged
coefficients are validated against the reality condition
\eqref{eq:real-check} and published atomically; the symbolic layer imports
them and proceeds to the next site. \emph{(An earlier version of this
back end delegated the same step to a Julia-based homotopy-continuation
package; several internal option and file names still reflect that
earlier design, though the currently used back end is Bertini only.)}

\textbf{Orchestration and batch execution.} A thin orchestration layer
isolates every numerical global used by the symbolic/asymptotic
construction on a per-task basis, allocates a dedicated working directory
and log files per tableau, and records detailed per-stage timing. A
shell-level batch launcher runs many tableaux (optionally all standard
Young tableaux of a given shape) as independent, parallel operating-system
processes, consolidating their individual results into a single aggregate
table under a file lock, atomically, so that a batch's partial results
remain safely readable while it is still running; a native desktop
application provides the same set of operations (single-tableau runs,
batch runs, the flip cross-check of Section~\ref{sec:validation-flip}, and
result visualization) through a graphical interface, without introducing
any code path distinct from the command-line tools.

\textbf{Visualization.} Bethe roots are visualized either in a notebook
(for interactive inspection and for the negated-overlay comparison used in
the flip check) or by a small, dependency-free tool that renders each
nesting level of \eqref{eq:bethe-roots} as a panel of a hand-written
vector-graphics document, automatically switching to a density-aggregated
rendering when the number of roots in a batch result is too large to plot
individually.

\begin{remark}[Why split the pipeline this way]
The division between a symbolic layer and a numerical back end mirrors the
division between Sections~\ref{sec:qsystem}--\ref{sec:asymptotics} (which
never involve floating-point arithmetic --- everything is exact
combinatorics and exact or symbolic algebra) and
Section~\ref{sec:homotopy} (the one place genuine numerical
approximation enters). Consequently, the only interface the two layers
need is the plain-text polynomial table described above: the symbolic
layer never needs to know \emph{how} a step's solution was obtained, only
that it converged, which is precisely what allowed the back end to be
replaced (Section~\ref{sec:implementation}'s remark on the earlier Julia
back end) without touching the mathematics of
Sections~\ref{sec:qsystem}--\ref{sec:asymptotics} at all.
\end{remark}


%=============================================================================
\section{Validation and Numerical Results}
\label{sec:results}
%=============================================================================

\subsection{Regression testing}

The numeric-string handling between Bertini's long-decimal output and the
symbolic layer's arbitrary-precision numbers is directly tested: a
100-digit decimal value is required to round-trip through export and
import with every digit intact, standard and scientific notations are
required to parse correctly, and malformed or non-finite input is
required to be rejected rather than silently coerced. Separately, a suite
of saved reference tableaux is rerun from scratch and compared against
their saved solutions \emph{level by level} (per Eq.~\eqref{eq:bethe-roots},
never flattened across nesting levels), using an exhaustive one-to-one
matching over root permutations at a default tolerance of $10^{-8}$.

\subsection{Validation of Conjecture~\ref{conj:flip}: the flip check}
\label{sec:validation-flip}

The empirical symmetry of Conjecture~\ref{conj:flip} is checked
numerically by computing, for a tableau and its flip, the
\textbf{symmetric Hausdorff distance} between one root set and the
negation of the other,
\begin{equation}
  \label{eq:hausdorff}
  d_H(A,B) \;=\; \max\Bigl\{\, \max_{a\in A}\min_{b\in B}|a-b| ,\;\; \max_{b\in B}\min_{a\in A}|a-b| \,\Bigr\},
  \qquad B=-\operatorname{Roots}[\mathrm{Flip}(\mathrm{SYT})],
\end{equation}
computed at high (100-digit) arithmetic precision, with a default pass
tolerance of $10^{-5}$; the check is applied automatically, at the level
of a whole aggregate batch of tableaux, and any pair exceeding tolerance
can be flagged and removed as a likely branch-jump artifact
(Section~\ref{sec:branch-jump-formal}). This is exactly the automated,
large-scale counterpart of the by-hand check carried out in
Section~\ref{sec:worked-example}: there, $d_H$ applied to
$\{+\tfrac{1}{2\sqrt3}\}$ and $-\{-\tfrac{1}{2\sqrt3}\}$ vanishes
identically, since the two single-element sets are literally equal;
for shapes too large to solve in closed form, the same comparison is
instead carried out numerically, at whatever precision the continuation
was run at, on every flip-related pair in a batch.

\begin{table}[H]
\centering
\begin{tabular}{@{}lll@{}}
\toprule
Young diagram & Tableau & Max.\ root difference \\
\midrule
$\{3,2,1\}$ & $\{\{1,2,6\},\{3,4\},\{5\}\}$ & $4.7\times10^{-17}$ \\
$\{3,2,1\}$ & $\{\{1,2,6\},\{3,5\},\{4\}\}$ & $2.9\times10^{-17}$ \\
$\{4,2,1\}$ & $\{\{1,2,3,4\},\{5,6\},\{7\}\}$ & $1.4\times10^{-14}$ \\
\bottomrule
\end{tabular}
\caption{Maximum one-to-one root differences (level by level) between freshly recomputed and saved reference solutions, all well below the $10^{-8}$ pass tolerance used for the comparison.}
\label{tab:accuracy}
\end{table}

Table~\ref{tab:accuracy} reports a representative set of from-scratch
reruns compared against saved reference solutions; all three cases agree
to well within tolerance. This establishes agreement with previously
computed (and themselves once machine-precision-validated) references to
the stated tolerance; it does not, by itself, certify absolute accuracy at
the full working precision used internally.

\subsection{Performance profile of a single tableau}

Table~\ref{tab:profile} reports a non-overlapping timing breakdown for one
complete computation of a seven-box tableau, totaling $8.95\,\mathrm{s}$,
with recovered roots matching a saved reference to within
$6.6\times10^{-17}$.

\begin{table}[H]
\centering
\begin{tabular}{@{}lrr@{}}
\toprule
Category & Seconds & Share \\
\midrule
Symbolic (Wolfram Language) construction \& equation preparation & 3.43 & 38.3\% \\
Bertini path tracking & 0.92 & 10.3\% \\
Other back-end process overhead & 1.04 & 11.6\% \\
Kernel launch/load/finalization overhead & 3.57 & 39.9\% \\
\bottomrule
\end{tabular}
\caption{Timing breakdown for one small tableau. Bertini's own path-tracking cost is a modest fraction of the total; fixed per-call kernel and symbolic-preparation overhead dominates at this problem size.}
\label{tab:profile}
\end{table}

For this particular (small) example, fixed overhead --- Wolfram Language
kernel initialization and one-time symbolic-preparation cost --- dominates
the wall-clock time, while the genuine numerical continuation cost
(Bertini path tracking) is comparatively small; this conclusion is
specific to the profiled example and has not been established for
substantially larger Young diagrams, where the polynomial degrees, and
hence the intrinsic path-tracking cost, grow.


%=============================================================================
\section{Discussion}
\label{sec:discussion}
%=============================================================================

\subsection{Summary}

We have given a formal account of the Q-system/Wronskian construction
used to solve the (supersymmetric) nested Bethe ansatz equations for
eigenstates labeled by standard Young tableaux: Q-functions built as
Casoratian determinants of a shared linear difference problem obey
conservation-law-type bilinear identities under the shift $u\to u\pm i$
(Eq.~\eqref{eq:wronskian-conservation}), the discrete analogue of which,
for the graded (bosonic/fermionic) elementary blocks of
Section~\ref{sec:susy}, generates the full Q-system by the Wronskian
recursion \eqref{eq:reduction-bos}--\eqref{eq:reduction-ferm}; the marked-point
freedom in splitting elementary building blocks into bosonic and
fermionic families corresponds to a choice of Kac--Dynkin grading; a
single chain inhomogeneity, continued from an asymptotically exact
combinatorial solution down to the physical homogeneous point, converts
the remaining numerical work into a sequence of parameter homotopy
continuations, formalized through the Davidenko equation and realized
numerically by Bertini's adaptive multiple-precision tracker; and an
independent rigged-configuration/string-hypothesis combinatorics supplies
both an approximate analytic cross-check and an exact combinatorial
symmetry (the flip involution) used as an empirical validation criterion.

\subsection{Limitations}

The physical identity of the bosonic/fermionic split (Remark following
Eq.~\ref{eq:susywronskian}) could not be settled from the codebase alone.
The flip symmetry of Conjecture~\ref{conj:flip} is used, and validated
numerically, but is not derived here from the representation theory of
the underlying algebra. The branch-jump failure mode of
Section~\ref{sec:branch-jump-formal} is mitigated by more conservative
continuation settings but is not eliminated by any formal guarantee, and
should continue to be screened for via the flip check on any newly
computed batch of tableaux. Performance characterization to date covers
only small tableaux; whether fixed overhead continues to dominate genuine
path-tracking cost at large Young-diagram size remains open.

\subsection{Outlook}

Two directions follow directly from the results above. First, since the
dominant cost observed in Table~\ref{tab:profile} is fixed per-call
overhead rather than path-tracking time, reusing a single initialized
symbolic-layer worker across many tableaux (rather than paying kernel
startup once per tableau) is a natural, low-risk optimization for batch
computations. Second, extending the performance profile of
Section~\ref{sec:results} to substantially larger Young diagrams would
determine whether the numerical continuation itself, rather than fixed
overhead, becomes the dominant cost once the Q-system's polynomial
degrees grow, which would in turn indicate where further numerical
effort (tighter path-tracking control, higher base precision, or an
adaptive step-size schedule tuned specifically to this problem's
observed branch-jump risk) is best directed.

\end{document}
